\documentclass[11pt]{article}
\usepackage[T1]{fontenc}
\usepackage[utf8]{inputenc}
\usepackage{amsmath,amssymb,amsthm,mathtools}
\usepackage[margin=1in]{geometry}
\usepackage{booktabs}
\usepackage{microtype}
\usepackage{enumitem}
\usepackage[font=small,labelfont=bf]{caption}
\usepackage{xcolor}
\usepackage[obeyspaces,spaces]{url}
\usepackage{hyperref}
\hypersetup{colorlinks=true, linkcolor=black, citecolor=black, urlcolor=blue}
\usepackage[capitalise,noabbrev]{cleveref}
\setlength{\emergencystretch}{2.5em}

\newtheorem{theorem}{Theorem}
\newtheorem{proposition}[theorem]{Proposition}
\newtheorem{lemma}[theorem]{Lemma}
\theoremstyle{definition}
\newtheorem{definition}[theorem]{Definition}
\newtheorem{remark}[theorem]{Remark}

\newcommand{\Z}{\mathbb{Z}}
\newcommand{\Q}{\mathbb{Q}}
% Facts not yet available (the DOI).
\newcommand{\pending}[1]{\textcolor{red}{#1}}

\title{Imaginary quadratic fields of class number at most 3500:\\ an independent extension, toward kernel verification}
\author{Peter W. Murphy}
\date{November 2026 (v2.0.0)\\[2pt] \small Repository \url{https://github.com/DrMurphyIsIn/boundaries-of-class}, directory \texttt{class-number-3500/}\\ \small DOI \url{https://doi.org/10.5281/zenodo.23924566} (v2.0.0)\\ \small (all versions: \url{https://doi.org/10.5281/zenodo.23550257})}

\begin{document}
\maketitle

\begin{abstract}
Assume the explicit bound for real zeros of Dirichlet $L$-functions in the second part of Theorem~1 of Kimura's preprint
\cite{Kimura2026}, which its author reports as checked in Lean by two independent kernels and which is not
refereed. Then exactly $50\,299\,869$
imaginary quadratic fields have class number at most $3500$, and the largest absolute discriminant among them is
$3\,495\,584\,467$, of class number $3477$. This extends Jaswal's classification of the fields with $h \le 3000$
\cite{Jaswal2026} by his own method and on the same footing, and it reproduces his list and Kimura's list for
$h \le 2000$ exactly. Every computational step of the ranges not covered by a cited result was re-implemented independently,
sharing no code with Kimura's or Jaswal's programs, and each was cross-checked against a second program or against the
published data. The same conclusion follows without the sieve from the real case of OpenAI's zero-free half-plane
$\operatorname{Re} s > 7/8$, whose Lean proof we have replayed with two kernels (its manuscript is not refereed); on
that route the analytic bound for large discriminants is kernel-checked in a companion paper. We release the
programs, the logs and the full list, together with a precise account of which steps are kernel-checked and which are not. The aim is to give the
community a basis from which this classification can be carried to a fully kernel-checked one.
\end{abstract}

\section{Introduction}

In October 2026 several arguments excluding Landau--Siegel zeros appeared, beginning with one released by OpenAI
\cite{OpenAISZ}; Terence Tao sketched the proof \cite{TaoSketch}. He then posed a challenge \cite{TaoChallenge}: for each
of these arguments, to use an AI tool to work out an effective upper bound for a hypothetical tenth imaginary quadratic
discriminant of class number one, without Heegner's methods, as a way of exploring the new arguments. Any amount of AI use
was permitted provided it was disclosed, and Tao asked contributors not to prioritise speed over quality: ``this is not
intended to be a race regarding who is `first' to claim any given bound.'' He also noted that a low enough threshold could
give a new solution of the class number one problem, possibly one formalizable in Lean.

The present paper does not answer that challenge directly. It belongs to a line of work that grew out of the same
explicit Siegel-zero bounds: once the zero bound is explicit, it bounds the discriminants of fields with any given class
number, and classifications far beyond the previously known range become a finite computation. We have tried to write it
in the spirit Tao asked for, with full attribution, a disclosed account of the AI's role, and an exact statement of what
has and has not been machine-checked.

It is the sixth entry of a repository, \emph{The Boundaries of Class}, that collects our work on this question. The first
three entries are kernel-checked proofs of the class number one theorem, by two routes \cite{MurphyFirst,MurphyCN1}, and
of the list of fields of class number two \cite{MurphyHCN}; the fourth is a survey of the bounds posted in the challenge
\cite{MurphyComparison}; the fifth proves explicit two-sided bounds for the class number of every imaginary quadratic
field, kernel-checked \cite{MurphyBounds}. In those entries the kernel checks the whole proof. Here it cannot, and much of
this paper is about saying exactly where the result stands.

Two recent preprints set the frontier. Kimura \cite{Kimura2026} made explicit, and checked in Lean with two independent kernels
according to his account, the argument of \cite{OpenAISZ}, using its Lean formalization \cite{OpenAIRepo}. He then determined, unconditionally, all
imaginary quadratic fields with $h \le 2000$: there are $16\,453\,493$, the largest $|D|$ being $1\,062\,350\,083$.
Jaswal \cite{Jaswal2026} assumed the same zero bound, which we call \emph{(K)} below, and extended the classification to
$h \le 3000$: there are $36\,993\,853$, the largest $|D|$ being $2\,562\,211\,723$. Both rest on Watkins's sieve
\cite{Watkins2004} and on the unconditional tabulation of Mosunov and Jacobson \cite{MJ16}. Earlier, on 8 October, the repository zeta-lab
had posted the fields with $h \le 1500$ under the full $7/8$ half-plane \cite{ZetaLab2026}; Jaswal's lists agree with it.

We extend Jaswal's classification to $h \le 3500$. The method is his; the footing is Kimura's (K); the tabulation is
Mosunov and Jacobson's. What we add is threefold:
\begin{itemize}[nosep]
\item an \emph{independent implementation} of every computational step of ranges (1), (3) and (4) below, sharing no
code with Kimura's or Jaswal's (range~(2) and the footing (K) are cited, not re-derived);
\item a \emph{cross-check of every step} against a second program or against the published data, including an
exact reproduction of both earlier lists and of Jaswal's sieve counts; and
\item a \emph{statement of status}, step by step, of what is kernel-checked, what is certified by interval arithmetic, and what
rests on agreement between programs, offered as a map for a fully kernel-checked version.
\end{itemize}

\begin{theorem}\label{thm:main}
Assume (K): for every primitive real Dirichlet character $\chi$ of conductor $q \ge 3$ and every real zero $\beta$ of
$L(s,\chi)$, if $\log q \ge 350131/10000$ then $(1-\beta)\log q \ge 1193/200000$ \cite[Theorem~1]{Kimura2026}. Then
exactly $50\,299\,869$ imaginary quadratic fields have class number at most $3500$. The largest absolute
discriminant among them is $3\,495\,584\,467$, with class number $3477$.
\end{theorem}

The sub-lists for $h \le 3000$ and $h \le 2000$ coincide exactly with Jaswal's and with Kimura's (\cref{sec:checks}).
Under (K) the three numbers rest on all four ranges: range~(1) computes them, and range~(3), at bound $3500$
(\cref{sec:range3}), leaves no discriminant unsettled. Under the hypothesis (H) of Remark~\ref{rem:H} they do not depend
on range~(3).

\section{The cover}\label{sec:cover}

Write $n = |D|$ for a fundamental discriminant $D < 0$. Following Jaswal \cite[\S2--3]{Jaswal2026}, the theorem is proved
by covering $n \ge 3$ with four ranges, each closed by a different argument (\cref{tab:cover}).

\begin{table}[h]
\centering\small
\begin{tabular}{@{}lll@{}}
\toprule
range & $n = |D|$ & argument \\
\midrule
(1) & $3 \le n < 19\,800\,000\,000$ & direct computation of every $h(D)$ (two independent programs)\\
(2) & $19\,800\,000\,000 \le n < 2^{40}$ & Mosunov--Jacobson \cite{MJ16}: $\min L(1,\chi) \approx 0.17070$, so $h > 7640$\\
(3) & $2^{40} \le n \le 1\,606\,926\,912\,624\,720$ & a Watkins sieve at bound $N = 3500$ (Jaswal's procedure)\\
(4) & $n \ge 1\,606\,926\,912\,624\,721$ & (K) and Hecke's integral representation: $h \ge 3562$\\
\bottomrule
\end{tabular}
\caption{The four ranges. The boundary between (3) and (4) is $\lfloor e^{35.0131}\rfloor$, checked in interval
arithmetic.}\label{tab:cover}
\end{table}

\paragraph{Range (4).} This is Jaswal's Proposition~2.7 \cite{Jaswal2026}: from (K) and Hecke's representation of the
Dedekind zeta function, $h(D) \ge 3562$ whenever $\log|D| \ge 35.0131$. We re-derived the argument line by line (theta
inversion, the representation count, the Mellin split) and re-certified the numerical constant in interval arithmetic;
our enclosure, $g(e^{35.0131}) \in [3561.29778, 3561.29779]$, lies inside Jaswal's printed enclosure
$[3561.2977, 3561.2978]$.

\paragraph{Range (2).} Mosunov and Jacobson tabulated the class groups of all negative discriminants with $|D| < 2^{40}$
unconditionally \cite{MJ16}. They record the minimum of $L(1,\chi)$ over their whole tabulation, about $0.17070$; it is attained
at $D = -107\,415\,709\,003$, which lies in range~(2). By the class number formula it gives $h > 7640$ for
$19\,800\,000\,000 \le |D| < 2^{40}$. We cite this, as Kimura and Jaswal do, after reading it at source.

\paragraph{Range (1).} For every fundamental $D$ with $|D| < 19\,800\,000\,000$ we decide whether $h(D) \le 3500$, and if so
compute $h(D)$ exactly. The first program uses Jaswal's procedure \cite[\S3.1]{Jaswal2026}: a class-order test
(baby-step giant-step on the classes of small split primes, with the genus-theory refinement) certifies $h > 3500$ for most
$D$, and the reduced-form count gives $h(D)$ exactly for the rest. The second program is PARI/GP's \texttt{qfbclassno}
\cite{PARI}, whose manual states that its results are unconditionally correct for $|D| < 2\cdot10^{10}$.

\paragraph{Range (3).}\label{sec:range3} We sieve every $n$ in the range with $n \bmod 16 \in \{3,7,11,15,4,8\}$ and
$9 \nmid n$ (a superset of the fundamental ones) through Jaswal's three steps \cite[\S3.2]{Jaswal2026}, in his eleven
intervals with his parameters: (i) a class-level bound by residues modulo a primorial-type modulus, (ii) a per-$n$ bound
from split primes just above the sieving range, and (iii) a partial reduced-form count. Our implementation enumerates the
residue classes of step~(i) by a depth-first search that certifies whole subtrees by a lower bound, and so avoids
enumerating most classes. In every interval the number of $n$ treated must equal an independent direct count, and no $n$
may be left over. In all eleven intervals both conditions hold. Step~(iii) was reached by $920\,101$ values of $n$ in
$[2^{40}, 2^{41})$, by $12\,528$ in $[2^{41}, 2^{42})$ and by $5$ in $[2^{42}, 2^{43})$, and by none above, and it settled
every one of them.

\begin{remark}[A shorter route under the 7/8 hypothesis]\label{rem:H}
Let (H) be the statement that $L(\sigma,\chi) \ne 0$ for $7/8 < \sigma < 1$ and every real primitive Dirichlet
character $\chi$. It is the real case of the half-plane discussed in \cref{sec:kernel}, and it implies (K). Under (H),
ranges (3) and (4) are not needed, as Jaswal observes \cite[Remark~4.2]{Jaswal2026}. By his Corollary~2.3 and the
monotonicity in his Proposition~2.2, $h(D) \ge B(2^{40}, 7/8)$ for every fundamental discriminant with $|D| \ge 2^{40}$,
where
\[
  B(x,\beta) = 2\beta(1-\beta)\sum_{m\ge1}\Bigl(\alpha_m^{-\beta}\,\Gamma(\beta,\alpha_m)
  + \alpha_m^{\beta-1}\,\Gamma(1-\beta,\alpha_m)\Bigr), \qquad \alpha_m = 2\pi m^2/\sqrt{x}.
\]
Every term is positive, so a partial sum is a lower bound. We evaluated the first $6000$ terms in interval arithmetic
and bounded the tail, obtaining $B(2^{40}, 7/8) = 17\,357.92029\ldots$, in agreement with Jaswal's value and far above
$3500$. So, under (H), the theorem follows from range~(1), range~(2), the analytic step $h(D) \ge B(|D|,\beta)$ (Jaswal's
Proposition~2.2, from Hecke's integral representation, with his Corollary~2.3) and this one constant. Indeed, as Jaswal
notes, under (H) the tabulation of Mosunov and Jacobson alone determines every field with $h < B(2^{40},7/8)$, so that
range~(1) becomes our independent computation of the same list below $19\,800\,000\,000$. On this route the numbers in the
theorem do not depend on range~(3) at all, and the footing is a statement whose Lean proof we have kernel-checked
ourselves. It is the natural route toward a fully kernel-checked version, and its analytic half is already kernel-checked by a
different route: the fifth entry \cite{MurphyBounds} proves in Lean, from the same $7/8$ theorem and checked by two
kernels, that $h(D) \ge 17\,450$ for every imaginary quadratic field with $|D| \ge 2^{40}$. Jaswal's Proposition~2.2,
which we re-derived by hand, is then not needed on this route. What would remain to formalize is range~(1) and the
tabulation of range~(2).
\end{remark}

\section{Checks}\label{sec:checks}

\begin{enumerate}[leftmargin=*]
\item \emph{Range (1), two programs.} The two programs' lists of $(n, h)$ with $h \le 3500$ are byte-identical after
normalisation: $50\,299\,869$ lines. Both tile $[3, 19\,800\,000\,000)$ in $1980$ segments with no gaps. The count of
fundamental discriminants treated, $6\,018\,478\,352$, equals Jaswal's.
\item \emph{Reproduction of the earlier lists.} Our sub-list for $h \le 3000$ has the same SHA-256 digest as Jaswal's
published list ($36\,993\,853$ fields, largest $2\,562\,211\,723$), and our sub-list for $h \le 2000$ has the same digest as
Kimura's ($16\,453\,493$ fields, largest $1\,062\,350\,083$). Every per-$h$ count agrees.
\item \emph{Range (3), reproduction at $N = 3000$.} Before running at $3500$, we ran our sieve at Jaswal's bound $3000$ over
all eleven intervals. In every interval the number of $n$ treated and the number reaching step~(iii) equal Jaswal's
published values, and the totals equal his $535\,275\,800\,332\,314$. Repeated runs of two of the intervals gave byte-identical output.
\item \emph{Range (3), coverage at $N = 3500$.} A separate checker verifies that the intervals are consecutive integer ranges,
that they meet the exact (3)/(4) boundary, that each was run once at $N = 3500$, that $n$ treated equals the direct count,
and that no $n$ is left unsettled. The check passes.
\item \emph{Separate review.} The written plan, which fixed every check above and the outcomes that would stop the work
in advance, the statements implemented, and each completed range were reviewed by a separate AI session (same model family, same operator; not a human
referee), which re-ran the cover check itself and obtained a byte-identical result. It also re-ran the first interval
of range~(3) at $N = 3500$ from a copy of the pinned program, and its output was byte-identical to the run of record.
\end{enumerate}

\section{Toward a kernel-checked classification}\label{sec:kernel}

The result above is \emph{not} kernel-checked as a whole, and we want to be precise about where it stands.
\begin{itemize}[nosep]
\item \textbf{(K)} is Lean-checked, with two independent kernels, according to Kimura \cite{Kimura2026}; we have
not replayed his formalization. But (K) also follows, as Jaswal notes \cite{Jaswal2026}, from the zero-free half-plane
$\operatorname{Re} s > 7/8$ for Dirichlet $L$-functions announced by OpenAI \cite{OpenAIQRH}: a real zero then has
$\beta \le 7/8$, so $(1-\beta)\log q \ge (\log q)/8 > 4$ when $\log q \ge 35.0131$. That half-plane we have checked
ourselves. OpenAI's Lean proof of it, the theorem
\path{OAI.DirichletCharacter.LFunction_ne_zero_of_seven_eighths_lt_re}, was replayed from openai/math at commit \texttt{adc7f124} through Comparator (run unsandboxed) and accepted by
Lean's kernel and by nanoda on 7 October 2026. It also lies inside a later export that five checkers accepted
(leanchecker, leanchecker-paranoid, nanoda, con-leche and con-ron; again not sandboxed, and lean4lean could not be used)
\cite{MurphyLchi}. A kernel check certifies the formal proof, not the manuscript: the manuscript is unrefereed, and the
match between the formal statement and the mathematical one was checked by reading it. (K) likewise follows from
Cureton and Thorner's half-plane $\operatorname{Re} s > 31/32$ for Dedekind zeta functions
\cite{CuretonThorner2026}, also announced with a Lean formalization, which we have not built.
\item \textbf{Range (4)} is a short analytic argument plus one numerical constant, which we certified in interval
arithmetic. It is a natural first target for formalization; Kimura's analogous Theorem~10 is already Lean-checked
according to his account.
\item \textbf{Range (2)} rests on a published, unconditional, but large tabulation \cite{MJ16}.
\item \textbf{Range (1)} rests on agreement between two independent programs (our own and PARI/GP), and with the
published lists, not on a kernel.
\item \textbf{Range (3)} at $N = 3500$ was run by our sieve alone. The evidence for it is that the same code reproduces
Jaswal's counts at $N = 3000$ exactly, interval by interval, together with the coverage checks at $N = 3500$; it is not
kernel-checked. Under (H) of Remark~\ref{rem:H} it is not needed.
\end{itemize}
Small class numbers can be kernel-checked today: in earlier work we gave kernel-checked proofs that the imaginary
quadratic fields of class number one are exactly the nine known ones, by two routes \cite{MurphyFirst,MurphyCN1}, and likewise
for class number at most two \cite{MurphyHCN}; the fifth entry kernel-checks explicit bounds for every field \cite{MurphyBounds}. Scaling that to $h \le 3500$ is a
different matter. A direct kernel replay of the range-(3) sieve with current tools would, by our measurements, take of the
order of hundreds of days on a single workstation, so we have not attempted it. We release the inputs, programs and logs
in a form intended to make such an effort, by us or by others, as straightforward as possible, and we would be glad to
collaborate on it.

\section{What we do and do not claim}\label{sec:claims}

We claim the classification of \cref{thm:main} under (K), with the checks of \cref{sec:checks}, and the same list under
(H) by the route of Remark~\ref{rem:H}. We do not claim the method: it is Jaswal's, on Kimura's footing and Watkins's
sieve, with Mosunov and Jacobson's tabulation, and we have followed it closely. We do not claim that the result is
kernel-checked, and we have said in \cref{sec:kernel} which parts are and which are not. The bound $3500$ was our
choice, set by the computing budget for range~(3). This paper says nothing about real quadratic fields, about orders that
are not maximal, or about the Riemann hypothesis.

Finally, this paper describes the computation and its checks. Where it states a count or a digest, the released files are
the authority.

\subsection*{The author's check}
The author read back the paper and the record of checks on 2 November 2026, and approved them. At that date the figures that depend
on range~(3) were still to be confirmed; the cover check, and the separate review of it, confirmed them on 3 November
2026.

\section*{Data and code}
The list of all $50\,299\,869$ fields with $h \le 3500$, the programs for ranges (1), (3) and (4), the run logs with
checksums, and the cover checker are in the directory \texttt{class-number-3500/} of the repository \url{https://github.com/DrMurphyIsIn/boundaries-of-class}; the full list is deposited with the release, DOI \url{https://doi.org/10.5281/zenodo.23924566}.

\section*{Acknowledgments and AI disclosure}
The challenge is Terence Tao's \cite{TaoChallenge}. The footing (K) and the reduction to a finite search are Iwao
Kimura's \cite{Kimura2026}, building on the argument released by OpenAI \cite{OpenAISZ} and its Lean formalization \cite{OpenAIRepo}; Banks
\cite{Banks2026}, Kerr \cite{Kerr2026} and Green \cite{Green2026} have given independent proofs of the Siegel-zero bound. The method, its parameters and the
classification up to $3000$ that we extend are Manvir Jaswal's \cite{Jaswal2026}; we followed his paper and his published
data closely and checked our work against them throughout. The sieve goes back to Watkins \cite{Watkins2004}, and range~(2)
is the tabulation of Mosunov and Jacobson \cite{MJ16}. We thank the PARI/GP developers.

The authors we build on state their own use of AI, which we record as they give it. Kimura states that the explicit
argument, the Lean files, the sieve program and most of the text were done with Claude Opus 5.5 (Anthropic) through
Claude Code, under his direction \cite[\S8]{Kimura2026}. Jaswal writes that his paper ``was produced by a multi-agent AI
research system that the author built and ran, using Anthropic Claude Opus 5.5: apart from the work of others that it
cites, the ideas, proofs, computations and text are due to the system'' \cite{Jaswal2026}.

The programs, the checks and this paper were produced with extensive use of an AI assistant (Claude, Anthropic) under the
direction of the author. A separate AI session, of the same model family and under the same operator, reviewed the plan,
the statements, each completed range and this paper, and re-ran the cover check with its own copy. That
review is not a substitute for a human referee. The mathematical route is due to the people credited above.
The account of status, the attributions and the account of how the work was checked are ours, and we have tried to say
how each was obtained.

\bibliographystyle{amsplain}
\bibliography{refs}

\end{document}
